\section{Introduction}
\label{sec:introduction}

A bank that holds an internal model for market risk capital reports a one-day value at risk (VaR) and, since the revision of the Basel standard, an expected shortfall (ES) every day, and the supervisor reads the record of the VaR against the realised losses through a count of exceedances over the last 250 trading days \citep{bcbs1996,bcbsmar32}. A model that understates the tail in a stress period shows the failure as a cluster of exceedances, and the standard's remedy is a larger capital multiplier applied after the fact. The validation of such a model asks three questions of the forecasts of each day: whether the long-run frequency of exceedances equals the tail level \citep{kupiec1995}, whether the exceedances arrive independently \citep{christoffersen1998}, and whether the ES that accompanies the VaR is large enough, scored jointly with it by a function for which the true pair is the unique minimiser \citep{fissler2016,pattonziegelchen2019}.

The estimators that produce the forecasts differ in how fast they react to a change in volatility. Historical simulation looks back a fixed window of returns \citep{pritsker2006,boudoukh1998}, volatility-filtered models update a conditional scale every day \citep{bollerslev1986,mcneilfrey2000,barone1999}, and quantile regressions let the quantile follow its own autoregression \citep{engle2004}. None of them knows in advance that a regime has changed. The experiment that motivates this manuscript is the eighteen-model ladder of Contraste, a product that validates the risk models of a bank: seventeen estimators of VaR and ES and one neural quantile network, run strictly forward on six windows of public daily series and scored by the same back-tests and the same score. In the financial crisis window of 2007 to 2009, none of the eighteen models passes the Kupiec and conditional coverage tests (both p-values at least 0.05) at the 1\% or the 2.5\% level, as the ladder runs without any recalibration (\cref{tab:ladder-context}). The reading a validator draws is that every model is wrong in the stress, which is correct, and unhelpful, because the validator also has to say what a bank can do about it without replacing the model.

The gap this manuscript addresses is the following. Online conformal methods repair the long-run coverage of a prediction interval under a distribution shift of unknown form, by moving the miscoverage level after each day by a fixed multiple of the surprise \citep{gibbs2021,gibbs2024,zaffran2022,angelopoulos2023}. Their application to VaR exists \citep{fantazzini2024,aich2025}, on large panels of crypto-assets and of financial series against intervals of returns. What is missing for a bank is a protocol that wraps a model that is already in production, leaves its estimator and its refit schedule untouched, keeps the ES that the model produces, and is judged by the instruments of the supervisor (the Kupiec and conditional coverage tests, the independence test, the Acerbi--Szekely statistic and the joint score of VaR and ES) on the stress windows where models fail, against the model it wraps and against other recalibrations run on the same forecasts. The recalibrated VaR of the supervisor's own models also needs a statement of when the repair cannot work.

The proposed method is the wrapper of \cref{sec:proposed-method}. It keeps a working tail level that starts at the target, asks the unchanged model for its VaR and ES at the working level, observes whether the day was an exceedance, and moves the level by a fixed step times the difference between the target and the exceedance indicator, inside a band. An exceedance makes the next VaR deeper, a quiet day lets it relax. The step, the band and the form of the ES are fixed before any run. The contributions are the following.
\begin{enumerate}
\item A deterministic identity for the clipped recursion that bounds the long-run coverage error by the width of the band, the step and the number of days the band was active (\cref{prop:identity}); the bound is vacuous for windows of the lengths tested here, which means that the repair measured below is an empirical result and not a consequence of the guarantee.
\item A proof that a guarantee on the frequency of exceedances carries no information on their independence, with a construction in which the frequency is exactly the target and the independence test rejects with certainty (\cref{prop:independence}).
\item A fixed-point analysis for a model that misjudges the scale of the data: the working level converges to the level at which the model's quantile equals the truth, the band limits the scale error that can be repaired (a ratio of 1.75 at the 1\% level for a Student-$t$ model with five degrees of freedom), and the band, acting at its lower edge, raises the hit rate above the target by a computable amount (\cref{prop:fixed-point}).
\item The test of the wrapper on ten models of the ladder in three stress windows, against the unwrapped models and against two online recalibrations that act on the finished VaR, with the results of a generator of known truth and of an ablation of each component.
\end{enumerate}

The main result is that the wrapped models pass the Kupiec and conditional coverage tests in \tenPassAci{} of \tenRuns{} model-level-window cells of the stress windows against \tenPassUnw{} of \tenRuns{} unwrapped, with no independence failure introduced (\tenIndFailAci{} failures wrapped against \tenIndFailUnw{} unwrapped), and with the mean FZ0 score changing by at most \tenWorstFz\% in any stress window. The prediction stated for the three models HS-250, GARCH-t and filtered historical simulation holds in all \hypCells{} of their model-level cells. The wrapper does not make the models right: the ES stays understated, the delta-normal model gains no stress window at the 1\% level, and two recalibrations that act on the finished VaR pass as many cells (\cmpRollTenPass{} of \cmpRollTenN{} for a rolling split-conformal scale) or more (\cmpTrkLargeTenPass{} for a quantile tracker with a larger step, which breaks the score tolerance in \cmpTrkLargeTenFzOver{} cells). \Cref{sec:related-work} places the method among the published ones, \cref{sec:problem} states the problem, \cref{sec:proposed-method} and \cref{sec:theory} give the method and its properties, \cref{sec:experimental-design} the experiments, and \cref{sec:results,sec:ablation} their results.
